\subsection*{§. 188.\\ 测地线与测地子流形的像。}

线元形式 (I*) 的 $ds^2$ 所对应的测地线微分方程很容易积分。事实上，应用 §. 156 的一般观察，可以立即求出方程
\[
 \Delta_1\theta=1,
\]
的一个解 $\theta$，其中含有 $n-1$ 个真正任意的常数。确实，若令
\[
 U=\frac{(y_1-a_1)^2+(y_2-a_2)^2+\cdots+(y_{n-1}-a_{n-1})^2+y_n^2}{y_n},
\]
其中 $a_1,a_2,\ldots,a_{n-1}$ 为任意常数，则上一节（公式 (10)）已见
\[
 \Delta_1U=U^2
\]
因而（§. 156）
\[
 \theta=\int\frac{dU}{\sqrt{\Delta_1U}}=\log U
\]
给出所求的 $\Delta_1\theta=1$ 的解。

根据 §. 156 所述，把同样数目的任意常数（记为 $b_1,b_2,\ldots,b_{n-1}$）分别等于 $\theta$ 对参数 $a_1,a_2,\ldots,a_{n-1}$ 的导数，即可得到测地线的有限形式方程，即
\[
\tag{\ensuremath{\alpha}}
 y_i-a_i=\frac{b_i}{2}y_nU\quad (i=1,2,\ldots,n-1).
\]
此外，测地线的弧长 $s$ 与 $\theta$ 只相差一个加法常数；于是
\[
\tag{\ensuremath{\beta}}
 U=Ce^s,
\]
其中 $C$ 是可任意选取的常数。借助上述公式，可以把 $y_1,y_2,\ldots,y_n$ 表示为 $s$ 的函数。将 $(\alpha)$ 平方后求和，再加上 $y_n^2$，得到
\[
 y_nU=\frac{b^2}{4}y_n^2U^2+y_n^2,
\]
其中为简便记
\[
 b^2=b_1^2+b_2^2+\cdots+b_{n-1}^2.
\]
由 $(\beta)$ 得
\[
 y_n=\frac{Ce^s}{\frac{b^2C^2}{4}e^{2s}+1}
\]
再令 $C=\frac2b$，便有
\[
 y_n=\frac1{b\cosh s}.
\]
于是 $(\alpha)$ 化为
\[
 y_i-a_i=\frac{b_i}{b^2}\frac{e^s}{\cosh s},
\]
即
\[
 y_i=\frac{b_i}{b^2}\tgh s+c_i,
\]
其中 $c_i$ 为新的常数。因此得到：\emph{伪球空间的测地线由下列方程表示；这些方程把测地线上动点的坐标表示为弧长的函数：}
\[
\tag{13}
\left\{
\begin{aligned}
 y_i&=\frac{b_i}{b^2}\tgh s+c_i\quad (i=1,2,\ldots,n-1),\\[.5em]
 y_n&=\frac1{b\cosh s},
\end{aligned}
\right.
\]
其中 $b_1,b_2,\ldots,b_{n-1}$、$c_1,c_2,\ldots,c_{n-1}$ 为任意常数，并且 $b^2=b_1^2+b_2^2+\cdots+b_{n-1}^2$。

这条线显然位于 $S_2$ 中，而该空间由 $n-2$ 个线性方程确定：
\[
 b_2(y_1-c_1)=b_1(y_2-c_2),
 \qquad
 b_3(y_1-c_1)=b_1(y_3-c_3),\ldots
\]
\[
 b_{n-1}(y_1-c_1)=b_1(y_{n-1}-c_{n-1}),
\]
这些方程各表示一个垂直于极限超平面的超平面；又有
\[
 \sum_{i=1}^{n-1}(y_i-c_i)^2+y_n^2=\frac1{b^2},
\]
故该线上的各点到点 $(c_1,c_2,\ldots,c_{n-1},0)$ 的距离相等，而该点属于这个 $S_2$。因此称这样的线为圆；它与极限超平面 $y_n=0$ 正交相交于两个点（对应于 $s=\pm\infty$）：
\[
 \left(c_1\pm\frac{b_1}{b^2},\ c_2\pm\frac{b_2}{b^2},\ldots,c_{n-1}\pm\frac{b_{n-1}}{b^2},\ 0\right).
\]
于是得到：\emph{伪球空间测地线的像是垂直于极限超平面的圆；反之，每一个这样的圆都是一条测地线的像。}

如上可见，这一表示把有关任意 $n$ 的性质推广开来；上一章曾研究 $n=2$ 的情形。重复该章的论证可知，伪球空间中的任意两点毫无例外地确定一条测地线；在我们的表示中，它们的实际测地距离由一个交比之比的对数度量：该交比由像圆上的两个像点与该圆同极限超平面的两个交点构成。

现在考虑一个超球面 $\Sigma_{n-1}$
\[
 (y_1-a_1)^2+(y_2-a_2)^2+\cdots+(y_{n-1}-a_{n-1})^2+y_n^2=r^2
\]
其中心位于极限超平面上。若以垂直于极限超平面的超平面截超球面 $\Sigma_{n-1}$，便得到一个球形子流形 $\Sigma_{n-2}$；其中心是从 $\Sigma_{n-1}$ 的中心向截超平面所作垂线的垂足，因而仍在极限超平面上。再截 $\Sigma_{n-2}$，所用超平面垂直于 $y_n=0$，便得到球形子流形 $\Sigma_{n-3}$，其中心在 $y_n=0$ 上；如此继续，最终得到中心在极限超平面上的普通球面 $\Sigma_2$。每个这样的球面都含有双无穷多个垂直于极限超平面的圆，因而相对于它的任一点都表示一个测地曲面。由此也显然可知，这些测地曲面的曲率为常数，且等于空间的曲率。同样，上述一切球形子流形
\[
 \Sigma_{n-1},\Sigma_{n-2},\ldots,\Sigma_2
\]
各表示一个测地子流形，即它们的一切测地线也是环境空间的测地线；此外，它们都具有等于空间曲率的常黎曼曲率。
